\documentclass[10pt,a4paper]{amsart}
\usepackage[margin=19mm]{geometry}
\usepackage{amsmath,amssymb,mathtools}
\usepackage[scr=rsfs,cal=euler]{mathalfa}
\usepackage{xfrac}
\usepackage[unicode,hidelinks,bookmarks=false]{hyperref}
\newcommand{\ie}{i.e.\ }
\newcommand{\cf}{cf.\ }
\newcommand{\resp}{respectively\ }
\newcommand{\Subsection}{\subsection}
\providecommand{\nobreakdash}{\nobreak\hbox{-}}
\setlength{\emergencystretch}{2em}
\multlinegap0pt
\usepackage{amssymb}
\newtheorem{thm}{Theorem}
\DeclareMathOperator{\Sp}{Sp}
\newcommand{\hh}{\tilde{h}}
\newcommand{\simc}{\!\!\sim}
\newcommand{\xb}{\mathbf{x}}
\title{On actions of tori and quaternionic tori on products of spheres}
\author{Anton Ayzenberg, Dmitry Gugnin}
\date{Source: arXiv:2309.05611v1 (CC BY 4.0)}
\begin{document}
\maketitle

\noindent\emph{Source and licence.} On actions of tori and quaternionic tori on products of spheres. Version arXiv:2309.05611v1 (CC BY 4.0), distributed by arXiv under the Creative Commons Attribution 4.0 International licence, http://creativecommons.org/licenses/by/4.0/, as recorded in the arXiv licence metadata for this exact version and shown on the abstract page https://arxiv.org/abs/2309.05611v1. Full licence notice: \texttt{licenses/arxiv-2309.05611-CC-BY-4.0.txt}.\par\smallskip

\noindent\textit{Editorial scope.} Counted unit: the article's thm thmJoins (source0.tex lines 321--329). The article's complete original proof of it is delivered with it (source0.tex lines 331--354, one \texttt{proof} environment of 24 lines). No mathematics is written, repaired or re-derived: the statement and the proof are the article's own lines, in the article's own notation, and no retained line is substituted or re-flowed.  Retained uncounted as the source-local context the proof needs: lines 182--187; lines 317--319.  Nothing was omitted from any retained span, so the package carries no omission record.  No retained line contains a citation, so the wrapper carries no bibliography and no bibliography entry.  No retained line contains a figure environment, an image-inclusion command or drawing code, so no image is delivered and the package declares no asset dependency.  Editorial formatting only, in addition: an AMS \texttt{amsart} preamble that restates the article's own declarations and macros used by the retained lines, namely the environments \texttt{thm} on separate counters, together with the standard packages those lines need.  The retained environments print in this document as Theorem 1 in order of appearance, whereas the article prints the same items with the numbering of its own full document.  The complete native source of the article as distributed in the arXiv e-print for this exact version is bundled unchanged as \texttt{sources/146853/source0.tex}.
\begin{thm}\label{thmQuaternionic}
The quotient space $S^{m_1}\times S^{m_2}\times \ldots \times S^{m_k}/\Sp(1)^{k-1}$ is homeomorphic to the sphere $S^m, m = m_1+\ldots + m_k - 3(k-1)$. The canonical projection to the orbit space is given by the formula:
\begin{equation}\label{eqTwoStars}
((\xb_1, q_1), (\xb_2, q_2), \ldots, (\xb_k, q_k)) \mapsto \frac{(\xb_1, \xb_2, \ldots, \xb_k, q_1q_2\ldots q_k)}{\sqrt{ |\xb_1|^2 + |\xb_2|^2 +  \ldots +  |\xb_k|^2 + |q_1q_2\ldots q_k|^2 }}.
\end{equation}
\end{thm}

\begin{equation}\label{eqCodiagonalAcn}
(x_1,1-t_1, q_1r_1^{-1},t_1; x_2,1-t_2, r_1q_2r_2^{-1},t_2; \ldots; x_k,1-t_k, r_{k-1}q_k,t_k).
\end{equation}

\begin{thm}\label{thmJoins}
The orbit space $\prod_{i=1}^{k}(X_i\ast G)/G^{k-1}$ is homeomorphic to the join
$X_1\ast X_2\ast \ldots\ast X_k\ast G$. The canonical projection onto the space of orbits is given by the formula:
\begin{multline}\label{eqJoinMap}
(x_1,1-t_1, q_1,t_1; x_2,1-t_2, q_2,t_2; \ldots; x_k,1-t_k, q_k,t_k) \mapsto \\
\mapsto (x_1,1-t_1;x_2,1-t_2;\ldots; x_k,1-t_k; q_1q_2\cdots q_k, t_1t_2\cdots t_k)/A,
\end{multline}
where $A=t_1\cdots t_k+\sum\nolimits_{i=1}^{k}(1-t_i)$ is the normalizing factor.
\end{thm}

\begin{proof}
Recall that the join of spaces $Y_1,\ldots,Y_m$ is the identification space
\[
Y_1\times\cdots\times Y_s\times \Delta^{m-1}/\simc,
\]
where $\Delta^{m-1}$ is the standard simplex with barycentric coordinates $(s_1,\ldots,s_m)$, $s_i\geqslant 0$, $\sum s_i=1$, and the equivalence relation $\sim$ is generated by the conditions
\[
(y_1,\ldots,y_l,\ldots,y_m,(s_1,\ldots,s_m))\sim (y_1,\ldots,y_l',\ldots,y_m,(s_1,\ldots,s_m)), \text{ if }s_l=0
\]
for any $l\in[m]$. Notice that $t_i\in[0;1]$ in~\eqref{eqJoinMap}, therefore $A>0$. Therefore, all coefficients $\frac{1-t_i}{A}$, $i\in[k]$ and $\frac{t_1\cdots t_k}{A}$ are nonnegative and sum to $1$ due to the choice of the normalizing factor $A$. Hence formula~\eqref{eqJoinMap} provides a well-defined continuous map of the form
\[
h\colon\prod\nolimits_{i=1}^{k}(X_i\times \Delta^1\times G)\to X_1\times\cdots\times X_k\times G\times \Delta^{k-1}.
\]
It is easily seen that the map $h$ is surjective. Let us check that $h$ descends to a well-defined map of the quotient spaces, the joins from the statement of the theorem. If $t_i=1$ for some $i\in[k]$, then the corresponding factor $X_i$ collapses to one point both in the space $X_i\ast G$ and in the space $X_i\ast\cdots \ast X_k\ast G$, so equivalent points of this kind are mapped to equivalent points. If $t_i=0$, then $t_1\cdots t_k=0$, and, similarly, equivalent points are mapped to equivalent ones. Therefore, the formula~\eqref{eqJoinMap} descends to a well-defined continuous map
\[
\hh\colon (X_1\ast G)\times (X_2\ast G)\times \ldots \times (X_k\ast G)/G^{k-1} \to X_1\ast X_2\ast \ldots\ast X_k\ast G.
\]
Since all the spaces appearing in this formula are Hausdorff compact, it suffices to show that the map $\hh$ is bijective. The surjectivity of $\hh$ follows from the surjectivity of~$h$.

Let us prove that $\hh$ is injective. Generally, the proof is similar to the proof of Theorem~\ref{thmQuaternionic}. Assume that
$\hh(x_1,t_1,q_1,\ldots,x_k,t_k,q_k) = \hh(x_1',t_1', q_1',\ldots, x_k',t_k',q_k')$. Looking at the real parameters $t_1,\ldots,t_k$ we see that formula~\eqref{eqJoinMap} defines a homeomorphism of the $k$-dimensional cube onto the $k$-dimensional simplex. Therefore the equalities $t_i=t_i'$ hold for all $i\in[k]$.

If $0<t_i<1$ for all $i$, then the assertion of injectivity reduces to the homeomorphism $G^k/G^{k-1}\cong G$ given by the formula $[(q_1,\ldots,q_k)]\mapsto q_1\cdots q_k$. If $t_i=1$ for some $i$, then the fiber $X_i$ collapses on both sides of the map. If $t_1\cdots t_k=0$, then $t_i=0$ for some $i\in[k]$. This means that the $i$-th component $G_i$ of the product $G^k$ collapses in the target of $\hh$. Taking the quotient of $G^k$ simultaneously by $G_i$ and the action of the subgroup $G^{k-1}$ described by the formula~\eqref{eqCodiagonalAcn}, we see that the entire group $G^k$ collapses into a point. Therefore, the identifications on the face of the simplex $\{t_1\cdots t_k=0\}$ are the same in the space $X_1\ast X_2\ast \ldots\ast X_k\ast G$ and in the space $(X_1\ast G)\times (X_2\ast G)\times \ldots \times (X_k\ast G)/G^{k-1}$. Injectivity is proved.
\end{proof}

\end{document}
